% =====================================================================
%  نسخه‌ی فارسی مقاله — باید با XeLaTeX کامپایل شود (نه pdfLaTeX).
%  در Overleaf: Menu -> Compiler -> XeLaTeX
%  این فایل به‌صورت محلی قابل کامپایل‌سنجی نبود (ابزار xelatex/xepersian
%  روی این سیستم نصب نیست)؛ لطفاً قبل از استفاده در Overleaf کامپایل و
%  بررسی شود.
% =====================================================================
\documentclass[a4paper,11pt]{article}

\usepackage{xepersian}
\settextfont{Vazirmatn}
\setlatintextfont{Latin Modern Roman}

\usepackage{amsmath}
\usepackage{amssymb}
\usepackage{algorithm}
\usepackage{algpseudocode}
\usepackage{graphicx}
\usepackage{booktabs}
\usepackage{tabularx}
\usepackage[table]{xcolor}
\usepackage{geometry}
\geometry{margin=2.4cm}
\usepackage[hyperfootnotes=false,colorlinks=true,linkcolor=blue!55!black,citecolor=blue!55!black,urlcolor=blue!55!black]{hyperref}

% اگر فایل تصویر موجود نبود، یک جعبه‌ی جانشین نشان بده تا سند همیشه کامپایل شود
\newcommand{\safefig}[2]{%
  \IfFileExists{#2}{\includegraphics[width=#1]{#2}}{%
    \fbox{\parbox[c][2.6cm][c]{\dimexpr#1-2\fboxsep-2\fboxrule\relax}{%
      \centering\ttfamily\detokenize{#2}\\[3pt]\normalfont(این تصویر را بارگذاری کنید)}}}%
}

\begin{document}

\begin{center}
{\Large\textbf{رتبه‌بندی منتقل می‌شود، آستانه نه:}}\\[4pt]
{\large تطبیق بدون‌نظارت در زمان آزمون برای آشکارسازی دیپ‌فیک صوتی میان‌پیکره‌ای}\\[14pt]
رهام ایزدی‌دوست\\
{\small دانشگاه مانیپال جیپور --- \lr{izadidoostroham@gmail.com}}\\[8pt]
دکتر سومیت سریواستاوا\\
{\small دانشگاه مانیپال جیپور --- \lr{sumit.srivastava@jaipur.manipal.edu}}\\[8pt]
شوتا شارما\\
{\small دانشگاه مانیپال جیپور --- \lr{Shweta.sharma1@jaipur.manipal.edu}}
\end{center}
\vspace{10pt}

\begin{abstract}
آشکارسازهای دیپ‌فیک صوتی روی داده‌ی هم‌پیکره (in-domain) دقتی نزدیک به کامل
گزارش می‌کنند، اما روی گفتار برگرفته از پیکره‌های ندیده---با شرایط ضبط،
مولّد، و زبان متفاوت---به‌شدت افت می‌کنند، و همین امر کاربرد واقعی را محدود
می‌سازد. ما یک مشاهده‌ی تجربی ساده اما پرپیامد ارائه می‌دهیم: میان پیکره‌ها،
\emph{رتبه‌بندی} (ROC-AUC) یک آشکارساز خودنظارتی فاین‌تیون‌شده تا حد زیادی
منتقل می‌شود، حتی وقتی \emph{آستانه‌ی} تصمیم آن کالیبره نمی‌ماند. این یعنی
پیش‌بینی‌های پراطمینان‌ترِ خودِ مدل روی صدای هدف بدون‌برچسب قابل‌اعتمادند، و
همین امر تطبیق بدون‌نظارت در \emph{زمان آزمون} (TTA) پایدار را ممکن می‌سازد:
ما روی برچسب‌های شبه پراطمینان (از دنباله‌های رتبه) خودآموزی می‌کنیم و
سازگاری میان کانال را اعمال می‌کنیم، در حالی که فقط لایه‌های LayerNorm و سر
طبقه‌بند تطبیق می‌یابند---بدون هیچ برچسب هدف و بدون داده‌ی منبع در زمان
آزمون. ارزیابی روی پروتکل «یک-پیکره-بیرون» با \emph{چهار} هدف نگه‌داشته‌شده
انجام می‌شود---ASVspoof2019، یک پیکره‌ی مبتنی بر LibriSpeech با گفتار
مصنوعی، In-the-Wild، و یک مجموعه‌ی دیپ‌فیک عربی چندزبانه---با یک پیکره‌ی
کمکی ۳۸‌زبانه در هر استخر منبع، سه بذر تصادفی برای هرکدام، و شش نقطه‌ی
مقایسه (Tent، بازکالیبراسیون فقط-BatchNorm، فقط-خودآموزی، یک ستون‌فقرات
غیرخودنظارتی از صفر، و یک پایه‌ی خصمانه‌ی دامنه در زمان آموزش). این روش نرخ
خطای برابر (EER) را در سه از چهار هدف بهبود می‌دهد (In-the-Wild
$12.57\!\rightarrow\!10.83\%$، عربی $22.59\!\rightarrow\!20.74\%$، ASVspoof2019
$4.76\!\rightarrow\!3.69\%$) و روی هدف چهارم (LibriSpeech-TTS،
$34.36\!\rightarrow\!34.38\%$) خنثی است، نه مضر، جایی که ظاهراً شکاف دامنه از
جنسی نیست که سازوکار مبتنی‌بر اطمینان ما بتواند از آن بهره ببرد. TTA مبتنی‌بر
کمینه‌سازی آنتروپی (Tent) روی دو از چهار هدف فرومی‌پاشد؛ یک ستون‌فقرات
سبک‌RawNet2 بدون پیش‌آموزش خودنظارتی روی هر چهار هدف فرومی‌پاشد؛ و یک پایه‌ی
خصمانه‌ی دامنه (DANN) که با سه برابر محاسبات روش ما آموزش دیده، در سه از چهار
هدف عملکردی ضعیف‌تر دارد. تحلیل فاصله‌ی $\mathcal{A}$ی جانشین نشان می‌دهد
جداپذیری دامنه پس از تطبیق \emph{کاهش نمی‌یابد}---که تبیین «تغییرناپذیری
دامنه» را رد می‌کند و در عوض از روایت رتبه‌بندی-کالیبراسیون ما پشتیبانی
می‌کند. یک تفکیک زبانی روی ۳۸ زبان MLAAD (دیده‌شده و ندیده) نشان می‌دهد
بازخوانی کلیپ‌های جعلی میان زبان‌های دیده و ندیده یکنواخت است (۹۳٫۹٪ در برابر
۹۳٫۳٪)، که نشان می‌دهد استواری چندزبانه‌ی آشکارساز مصنوع پوشش زبانی هنگام
آموزش نیست.
\end{abstract}

\textbf{کلیدواژه‌ها:} آشکارسازی دیپ‌فیک صوتی، ضدجعل، تطبیق در زمان آزمون،
تعمیم میان‌پیکره‌ای، یادگیری خودنظارتی، آموزش خصمانه‌ی دامنه، EER.

\section{مقدمه}
گفتار مصنوعی اکنون آن‌قدر واقع‌نما شده که تأیید هویت صوتی را تهدید و کلاه‌برداری
را تسهیل می‌کند، و همین امر انگیزه‌ی ساخت آشکارسازهای قابل‌اعتماد است. روی
معیارهای استاندارد، پیش‌سرهای خودنظارتی (SSL) مانند wav2vec\,2.0 و XLS-R،
همراه با پس‌سرهای سبک، به خطای بسیار پایینی می‌رسند~\cite{wav2vec2aug,xlsrsls,aasist}.
اما این اعداد عمدتاً \emph{هم‌پیکره}‌اند: آموزش و آزمون روی یک پیکره. وقتی
مدل روی صدایی از منبعی متفاوت---میکروفون، کدک، زبان، یا مولّد جدید---به کار
گرفته شود، دقت به‌شدت افت می‌کند؛ شکاف تعمیمی که به‌خوبی مستند شده
است~\cite{asdg,crossdomain,itw}.

ما از یک اندازه‌گیری تشخیصی آغاز می‌کنیم. با گرفتن یک آشکارساز XLS-R
فاین‌تیون‌شده و ارزیابی آن روی پیکره‌ای ندیده، مشاهده می‌کنیم که \textbf{ROC-AUC
بالا می‌ماند در حالی که دقت در یک آستانه‌ی ثابت فرومی‌پاشد}: مدل همچنان
کلیپ‌ها را درست \emph{رتبه‌بندی} می‌کند (واقعی زیر جعلی) اما نقطه‌ی عملیاتی‌اش
برای دامنه‌ی جدید کالیبره نیست. این نامتقارنی، کلیدی است که یک راه‌حل عملی
را می‌گشاید. چون رتبه‌بندی منتقل می‌شود، پیش‌بینی‌های پراطمینان‌ترِ مدل روی
صدای هدف بدون‌برچسب قابل‌اعتمادند، پس می‌توان از آن‌ها برای تطبیق مدل با
دامنه‌ی هدف \emph{در زمان آزمون}، بدون هیچ برچسبی، استفاده کرد.

نسخه‌ی اولیه‌ی این ایده، که روی دو هدفِ نگه‌داشته‌شده ارزیابی شد، بهبودهای
پیوسته نشان داد. مقاله‌ی حاضر آن ارزیابی را به‌طور چشمگیری گسترش می‌دهد: دو
هدف نگه‌داشته‌شده‌ی دیگر (از جمله هدفی که، همان‌طور که در ادامه گزارش
می‌کنیم، روش روی آن کمکی \emph{نمی‌کند})، یک پیکره‌ی کمکی ۳۸‌زبانه که در هر
استخر منبع گنجانده شده تا استواری چندزبانه‌ی واقعی را بیازماید، یک مطالعه‌ی
حذفی روی ستون‌فقرات غیر-SSL، یک پایه‌ی خصمانه‌ی دامنه در زمان آموزش که
به‌جای فقط آزمایشی‌بودن تا همگرایی آموزش دیده، و اندازه‌گیری مستقیم اینکه
آیا تطبیق جداپذیری دامنه را در فضای نهفته تغییر می‌دهد یا نه. ما نتیجه‌ی
منفی روی یک هدف را شاهدی مهم درباره‌ی \emph{دامنه‌ی کاربرد} روش می‌دانیم، نه
چیزی برای حذف: این نشان می‌دهد TTA مبتنی‌بر اطمینان، ویژگی خاصی از جابه‌جایی
منبع$\to$هدف را بهره می‌برد (رانش کالیبراسیون زیر یک رتبه‌بندی حفظ‌شده)، و
این ویژگی در همه‌ی جفت‌های پیکره یکسان برقرار نیست.

سهم‌های ما عبارت‌اند از:
\begin{itemize}
\item یک یافته‌ی تجربی---\emph{رتبه‌بندی میان‌پیکره‌ای منتقل می‌شود حتی وقتی
آستانه منتقل نمی‌شود}---و پیامد آن برای تطبیق بدون‌نظارت.
\item یک روش TTA ساده و پایدار: خودآموزی روی برچسب‌های شبه از دنباله‌های
پراطمینان به‌علاوه‌ی سازگاری کانال، با تطبیق فقط LayerNorm و سر.
\item ارزیابی چهارهدفه، چندزبانه، چندبذری با پروتکل یک-پیکره-بیرون با شش
نقطه‌ی مقایسه (فقط-منبع، Tent، بازکالیبراسیون فقط-BatchNorm، فقط-خودآموزی،
یک ستون‌فقرات غیر-SSL، و DANN) و یک کران بالای اوراکل، با گزارش صادقانه
شامل هدفی که روش روی آن کمکی نمی‌کند.
\item یک تحلیل سازوکاری---فاصله‌ی $\mathcal{A}$ی جانشین و یک کاوشگر منشأ
دامنه پیش و پس از تطبیق---که نشان می‌دهد بهبود از حذف هویت دامنه
\emph{نمی‌آید}، سازگار با روایت رتبه‌بندی-کالیبراسیون ما.
\item یک تفکیک زبانی روی ۳۸ زبان که بازخوانی یکنواخت کلیپ‌های جعلی را میان
زبان‌های دیده و ندیده در آموزش نشان می‌دهد.
\end{itemize}

\section{کارهای پیشین}
\textbf{ضدجعل مبتنی بر SSL.} پیش‌سرهای منجمد یا فاین‌تیون‌شده‌ی
wav2vec\,2.0 / XLS-R / WavLM با پس‌سرهای AASIST یا توجهی، بر معیارهای اخیر
مسلط‌اند~\cite{wav2vec2aug,xlsrsls,aasist,rawnet2}. این‌ها هم‌پیکره قوی‌اند
اما پیکره‌ویژه‌اند؛ ما از این خانواده (XLS-R فاین‌تیون‌شده) به‌عنوان
آشکارساز خود استفاده می‌کنیم و افزون بر آن می‌آزماییم که آیا شکاف
میان‌پیکره‌ای که با آن مواجه است اصلاً یک پرسش SSL است، با مقایسه در برابر
یک ستون‌فقرات غیر-SSL از صفر (بخش~\ref{sec:backbone}).

\textbf{تعمیم دامنه (DG) برای ADD.} روش‌های زمان‌آموزش، واقعی و جعلی را در
دامنه‌های منبع هم‌تراز/جدا می‌کنند (مثلاً تجمیع-جدایی~\cite{asdg}) یا
ویژگی‌های SSL را ترکیب می‌کنند~\cite{crossdomain}. نمونه‌ی متعارف این
خانواده، آموزش خصمانه‌ی دامنه با بازگشت گرادیان (DANN) است~\cite{dann}، که
یک تمایزگر دامنه را خصمانه در برابر استخراج‌کننده‌ی ویژگی آموزش می‌دهد تا
منبع و هدف در فضای نهفته خطی‌ناپذیر شوند. این روش‌ها شکاف را در زمان آموزش
هدف قرار می‌دهند، به داده‌ی هدف (بدون‌برچسب) \emph{در حین} آموزش نیاز
دارند نه در استقرار، و---همان‌طور که با یک پایه‌ی DANN آموزش‌دیده تا
همگرایی روی داده‌ی خودمان تأیید می‌کنیم (بخش~\ref{sec:dann})---به‌روشنی از
یک جایگزین بسیار ارزان‌تر در زمان آزمون بهتر نیستند.

\textbf{تطبیق در زمان آزمون.} TTA یک مدل آموزش‌دیده را بدون بازآموزی از صفر
با توزیع هدف بدون‌برچسب تطبیق می‌دهد. کمینه‌سازی آنتروپی روی پارامترهای
نرمال‌سازی (Tent)~\cite{tent} پرکاربردترین دستور در بینایی است اما،
همان‌طور که نشان می‌دهیم، برای آشکارسازی دیپ‌فیک صوتی ناپایدار است: بدون
کف برای اینکه مدل چقدر می‌تواند مطمئن شود، می‌تواند نمره‌ی تصمیم را به یک
ثابت واردشونده براند. بازکالیبراسیون آماره‌های دسته (رویکردهای
سبک-AdaBN)~\cite{adabn} آماره‌های نرمال‌سازی را روی داده‌ی هدف بدون هیچ
گام گرادیانی بازمحاسبه می‌کنند؛ ما یک حذفی منطبق \emph{فقط-BN} می‌گنجانیم و
می‌یابیم که برای معماری ما بی‌اثر است، چون LayerNorm یک ترنسفورمر
فاین‌تیون‌شده هیچ آماره‌ی دسته‌ای برای بازمحاسبه ندارد---دلیلی معماری، نه
فقط تجربی، برای اینکه چرا تطبیق گرادیانی مبتنی‌بر اطمینان اینجا لازم است.
TTA ویژه‌ی صدا برای آشکارسازی دیپ‌فیک نسبتاً کم‌کاویده است؛ نشان می‌دهیم
کمینه‌سازی آنتروپی ساده روی این وظیفه ناپایدار است و جایگزینی پایدار و
رتبه‌محور پیشنهاد می‌کنیم که فقط پارامترهای آفاین LayerNorm و سر را تطبیق
می‌دهد.

\textbf{واگرایی دامنه و مجموعه‌داده‌ها.} ما ساخت فاصله‌ی $\mathcal{A}$ی
جانشین را از نظریه‌ی تطبیق دامنه‌ی بن‌دیوید و همکاران~\cite{bendavid}
وام می‌گیریم، که نشان می‌دهند خطای دامنه‌ی هدف یک طبقه‌بند با خطای منبع آن
به‌علاوه‌ی جمله‌ای که میزان جداپذیری منبع از هدف توسط یک تمایزگر دامنه را
اندازه می‌گیرد کران‌دار است. ما از این برای آزمودن تجربی اینکه آیا تطبیق ما
این جداپذیری را کاهش می‌دهد استفاده می‌کنیم (بخش~\ref{sec:divergence}). برای
پیکره‌های ارزیابی از ASVspoof2019-LA، پیکره‌ی In-the-Wild~\cite{itw}، و
پیکره‌ی چندزبانه‌ی MLAAD هم برای استخر منبع کمکی چندزبانه و هم برای یک هدف
واقعاً چندزبانه (دیپ‌فیک‌های عربی نگه‌داشته‌شده) و تشخیص زبانی استفاده
می‌کنیم.

\section{روش}
\subsection{صورت‌بندی مسئله}
$x$ یک موج صوتی با برچسب $y\in\{0,1\}$ باشد ($0$ واقعی، $1$ جعلی). یک
آشکارساز $f_\theta$، $x$ را به ویژگی‌های فریمی نگاشت می‌دهد، آن‌ها را به
یک نهفتگی $e\in\mathbb{R}^{d}$ تجمیع می‌کند، و یک احتمال جعلی خروجی
می‌دهد
\begin{equation}
s_\theta(x)=\sigma\!\left(w^\top e+b\right)=P_\theta(y{=}1\mid x),\qquad \sigma(z)=\tfrac{1}{1+e^{-z}} .
\end{equation}
آموزش روی یک پیکره‌ی \emph{منبع} برچسب‌دار $\mathcal{D}_S=\{(x_i,y_i)\}$
آنتروپی متقاطع را کمینه می‌کند
\begin{equation}
\mathcal{L}_{\mathrm{CE}}(\theta;\mathcal{D})=-\frac{1}{|\mathcal{D}|}\!\!\sum_{(x,y)\in\mathcal{D}}\!\!
\big[y\log s_\theta(x)+(1{-}y)\log(1{-}s_\theta(x))\big].
\end{equation}
در زمان آزمون یک پیکره‌ی هدف \emph{بدون‌برچسب}
$\mathcal{D}_T=\{x_j\}_{j=1}^{N}$ از توزیعی جابه‌جاشده (پیکره، کانال، یا
زبان متفاوت) دریافت می‌کنیم؛ نه برچسب هدف و نه داده‌ی منبع در دسترس است.

\subsection{آشکارساز}
رمزگذار $g_\psi$ همان XLS-R-300M است که فقط چهار لایه‌ی بالایی ترنسفورمر
آن آموزش‌پذیرند. با ویژگی‌های فریمی $\{h_t\}_{t=1}^{T}$، تجمیع زمانی
توجهی نهفتگی را می‌سازد
\begin{equation}
\alpha_t=\frac{\exp(v^\top h_t)}{\sum_{t'}\exp(v^\top h_{t'})},\qquad
e=\sum_{t=1}^{T}\alpha_t\,h_t,
\end{equation}
و سپس یک تصویرسازی و سر خطی معادله‌ی (۱). فاین‌تیون کردن فقط لایه‌های
بالایی بیشتر شکاف هم‌پیکره را می‌بندد در حالی که روی یک GPU قابل‌آموزش
می‌ماند. نکته‌ی مهم برای طراحی تطبیق زیر، لایه‌های نرمال‌سازی XLS-R از نوع
\emph{LayerNorm}‌اند، که روی بُعد ویژگی به‌ازای هر نمونه نرمال می‌کند و
هیچ آماره‌ی دسته‌ی جاری ندارد---برخلاف BatchNorm، که آماره‌های جمعیتی
کانالی را انباشته می‌کند که یک دسته‌ی هدف می‌تواند به‌سادگی بازنویسی کند.
همین واقعیت معماری است که یک پایه‌ی بازکالیبراسیون آماره‌ی دسته
(بخش~\ref{sec:main-results}، \emph{فقط-BN}) را اینجا بی‌اثر می‌کند: هیچ
آماره‌ی دسته‌ای برای بازمحاسبه نیست، فقط پارامترهای آفاین $\{\gamma,\beta\}$
که به سیگنال گرادیانی برای حرکت نیاز دارند. پس خودآموزی مبتنی‌بر اطمینان
انتخابی دلبخواه نیست، بلکه کمینه سازوکاری است که اصلاً می‌تواند این لایه‌ی
نرمال‌سازی خاص را تطبیق دهد.

\subsection{شکاف رتبه‌بندی--کالیبراسیون}
برای آستانه‌ی $\tau$ تصمیم $\hat y_\tau(x)=\mathbf{1}[s_\theta(x)\ge\tau]$
است، که دقتش به $\tau$ بستگی دارد. کیفیت رتبه‌بندی، در مقابل، بدون آستانه
است:
\begin{equation}
\mathrm{AUC}=\Pr_{x^{+}\sim\mathrm{fake},\,x^{-}\sim\mathrm{real}}\!\big[s_\theta(x^{+})>s_\theta(x^{-})\big],
\end{equation}
که نسبت به هر بازپارامترسازی اکیداً افزایشی $s_\theta$ ناوردا است. تجربتاً،
زیر جابه‌جایی منبع$\rightarrow$هدف، $\mathrm{AUC}$ بسیار کمتر از دقت در یک
آستانه‌ی ثابت افت می‌کند (شکل~\ref{fig:scores}: یک تای نماینده که AUC بالای
$0.89$ می‌ماند در حالی که دقت در $\tau{=}0.5$ به میانه‌ی هفتادها می‌افتد).
\emph{رتبه‌بندی بیش از کالیبراسیون منتقل می‌شود.} پس ما به \emph{ترتیب}
نمره اعتماد می‌کنیم، نه مقدار مطلق آن: چندک‌های افراطی برچسب‌های شبه
دقیق‌تری نسبت به یک آستانه‌ی ثابت می‌دهند، حتی وقتی $s_\theta$ کالیبره
نباشد. همان‌طور که بخش~\ref{sec:main-results} نشان می‌دهد، این سازوکار
همگانی نیست: در یکی از چهار هدف ما (LibriSpeech-TTS) AUC منبع از پیش تنها
حدود $0.72$ است، به‌اندازه‌ی کافی نزدیک به شانس که حتی چندک‌های افراطی هم
مطمئناً درست نیستند، و کل پیش‌فرض روش ضعیف‌تر است---که دقیقاً همان‌جایی است
که تطبیق دیگر کمک نمی‌کند.

\begin{figure*}[t]
\centering
\safefig{0.86\textwidth}{fig_score_dist.png}
\caption{نمرات آشکارساز (محور لوجیت؛ نمرات نزدیک افراط‌ها متمرکز می‌شوند)
برای کلیپ‌های واقعی در برابر جعلی روی پیکره‌ی \emph{ندیده}‌ی In-the-Wild
(بذر نماینده از تحلیل مقدماتی). \textbf{چپ:} مدل منبع بسیاری از کلیپ‌های
واقعی را از آستانه‌ی $0.5$ عبور می‌دهد (دقت در $\tau{=}0.5$ به‌وضوح کمتر از
AUC) اما همچنان بیشتر واقعی‌ها را زیر جعلی‌ها رتبه‌بندی می‌کند.
\textbf{راست:} تطبیق بدون‌نظارت در زمان آزمون توده‌ی واقعی را دوباره به
سمت دیگر آستانه می‌راند، و هم دقت@۰٫۵ و هم AUC را بالا می‌برد.
جدول‌های~\ref{tab:main}--\ref{tab:gaindecomp} میانگین‌های کامل ۳بذر،
۴هدف مورد استفاده در همه‌ی ادعاهای اصلی این مقاله را گزارش می‌کنند.}
\label{fig:scores}
\end{figure*}

\subsection{تطبیق در زمان آزمون}
$Q_\alpha$ چندک-$\alpha$ی نمرات هدف $\{s_\theta(x_j)\}_{j=1}^{N}$ باشد.
در هر دوره‌ی تطبیق، برچسب‌های شبه پراطمینان را از دنباله‌های افراطی
(دوباره) تخصیص می‌دهیم،
\begin{equation}
\tilde y_j=
\begin{cases}
0, & s_\theta(x_j)\le Q_{q},\\
1, & s_\theta(x_j)\ge Q_{1-q},\\
\varnothing, & \text{در غیر این صورت (نادیده گرفته می‌شود)},
\end{cases}
\end{equation}
و $\mathcal{C}=\{j:\tilde y_j\neq\varnothing\}$. خودآموزی، آنتروپی متقاطع
روی این مجموعه‌ی مطمئن است،
\begin{equation}
\mathcal{L}_{\mathrm{ST}}=-\frac{1}{|\mathcal{C}|}\sum_{j\in\mathcal{C}}
\big[\tilde y_j\log s_\theta(x_j)+(1{-}\tilde y_j)\log(1{-}s_\theta(x_j))\big].
\end{equation}
با یک تقویت تصادفی کانال $a(\cdot)$ (بهره‌ی تصادفی $+$ نویز جمعی) و
$p_\theta(x)=[\,1{-}s_\theta(x),\,s_\theta(x)\,]$، یک جمله‌ی سازگاری
تغییرناپذیری پیش‌بینی نسبت به کانال را اعمال می‌کند،
\begin{equation}
\mathcal{L}_{\mathrm{cons}}=\frac{1}{N}\sum_{j=1}^{N}\big\|\,p_\theta(a(x_j))-\mathrm{sg}\!\left[p_\theta(x_j)\right]\big\|_2^2,
\end{equation}
که در آن $\mathrm{sg}[\cdot]$ توقف‌گرادیان است. هدف تطبیق چنین است
\begin{equation}
\min_{\{\gamma,\beta\},\,\phi}\ \ \mathcal{L}_{\mathrm{ST}}+\lambda\,\mathcal{L}_{\mathrm{cons}},
\qquad \lambda=0.3,\ q=0.3,
\end{equation}
که فقط پارامترهای آفاین LayerNorm $\{\gamma,\beta\}$ لایه‌های فاین‌تیون‌شده
و سر $\phi$ را بهینه می‌کند؛ باقی $\psi$ منجمد می‌ماند. محدود کردن مجموعه‌ی
آموزش‌پذیر و استفاده‌ی صرف از دنباله‌های مطمئن، از فروپاشی واردشونده‌ی
$s_\theta(\cdot)\!\to\!\text{ثابت}$ که کمینه‌سازی آنتروپی مستعد آن است
جلوگیری می‌کند، که تجربتاً در برابر Tent تأیید می‌کنیم
(بخش~\ref{sec:main-results}). با این حال، همان‌طور که
بخش~\ref{sec:gaindecomp} نشان می‌دهد، $\mathcal{L}_{\mathrm{cons}}$ همیشه
سودمند نیست: زیر شکاف دامنه‌ی بزرگ (عربی) نجات‌بخش تطبیق است اما روی تنها
هدفی که سیگنال رتبه‌بندی از پیش ضعیف است (LibriSpeech-TTS) زیان خالص است،
برهم‌کنشی که پیش‌بینی نکرده بودیم و آن را به‌صورت یک محدودیت صادقانه گزارش
می‌کنیم. الگوریتم~\ref{alg:tta} روند را خلاصه می‌کند.

\begin{algorithm}[t]
\caption{تطبیق در زمان آزمون بر منیفلد واقعی (به‌ازای هر پیکره‌ی هدف)}
\label{alg:tta}
\begin{algorithmic}[1]
\Require مدل منبع $f_\theta$؛ هدف بدون‌برچسب $\mathcal{D}_T=\{x_j\}_{j=1}^{N}$؛
چندک $q$؛ وزن $\lambda$؛ دوره‌ها $E$؛ تقویت $a(\cdot)$؛ گام $\eta$
\Ensure مدل تطبیق‌یافته $f_\theta$
\State لایه‌ی LayerNorm بالایی $\{\gamma,\beta\}$ و سر $\phi$ را آموزش‌پذیر کن؛ بقیه را منجمد کن
\For{$e=1,\dots,E$}
  \State $s_j \gets s_\theta(x_j)\ \ \forall j$ \Comment{کل مجموعه‌ی هدف را نمره‌دهی کن}
  \State $Q_q,\,Q_{1-q}\gets$ چندک‌های $\{s_j\}_{j=1}^{N}$
  \State $\tilde y_j \gets 0$ اگر $s_j\!\le\!Q_q$؛\ \ $1$ اگر $s_j\!\ge\!Q_{1-q}$؛\ \ وگرنه $\varnothing$
  \State $\mathcal{C}\gets\{j:\tilde y_j\neq\varnothing\}$ \Comment{دنباله‌های مطمئن}
  \For{زیردسته‌ی $\mathcal{B}\subset\mathcal{D}_T$}
    \State $\mathcal{L}_{\mathrm{ST}}\gets\mathrm{CE}$ روی $\mathcal{B}\cap\mathcal{C}$ با هدف‌های $\tilde y$ \Comment{معادله‌ی (۶)}
    \State $\mathcal{L}_{\mathrm{cons}}\gets\frac{1}{|\mathcal{B}|}\sum_{x\in\mathcal{B}}\|p_\theta(a(x))-\mathrm{sg}[p_\theta(x)]\|_2^2$
    \State $\theta\gets\theta-\eta\,\nabla_\theta(\mathcal{L}_{\mathrm{ST}}+\lambda\mathcal{L}_{\mathrm{cons}})$
  \EndFor
\EndFor
\end{algorithmic}
\end{algorithm}

\subsection{اندازه‌گیری واگرایی دامنه}
\label{sec:method-divergence}
برای آزمودن \emph{چرایی} کمک‌کننده بودن تطبیق---اینکه آیا با تمایزناپذیر
کردن منبع و هدف در فضای نهفته کار می‌کند، همان‌طور که روش‌های DG خصمانه
صریحاً برای آن بهینه می‌شوند---فاصله‌ی $\mathcal{A}$ی جانشین را به‌پیروی از
بن‌دیوید و همکاران~\cite{bendavid} محاسبه می‌کنیم. نهفتگی‌های منجمد $e$ را
برای ۷۵۰ کلیپ منبع و ۷۵۰ کلیپ هدف استخراج می‌کنیم، یک طبقه‌بند دامنه‌ی خطی
$h$ را برای پیش‌بینی منشأ منبع-در-برابر-هدف زیر اعتبارسنجی متقاطع ۵تایی
آموزش می‌دهیم، و تعریف می‌کنیم
\begin{equation}
\hat d_{\mathcal{A}} = 2\left(1-2\,\epsilon_h\right),
\end{equation}
که در آن $\epsilon_h$ خطای نگه‌داشته‌شده‌ی طبقه‌بند دامنه است. $\hat
d_{\mathcal{A}}$ی بزرگ‌تر یعنی منبع و هدف در فضای نهفته راحت‌تر از هم
قابل‌تشخیص‌اند. این را هم از نهفتگی‌های مدل منبع و هم از نهفتگی‌های مدل
تطبیق‌یافته روی همان کلیپ‌ها، به‌ازای هر هدف محاسبه می‌کنیم
(بخش~\ref{sec:divergence}).

\section{تنظیمات آزمایشی}
\textbf{پروتکل.} یک-پیکره-بیرون روی چهار هدف قابل‌EER: ASVspoof2019-LA، یک
پیکره‌ی واقعی/چندسامانه‌ی TTS مبتنی بر LibriSpeech (D2)، In-the-Wild~\cite{itw}،
و مجموعه‌ی دیپ‌فیک صوتی عربی (ArAD)، هدف چندزبانه‌ی ما. برای هر هدف
نگه‌داشته‌شده، استخر منبع \emph{سه پیکره‌ی دیگرِ} قابل‌EER را با یک زیرمجموعه‌ی
MLAAD ۳۸‌زبانه (فقط جعلی، آب‌پرکرده میان زبان‌ها) که در هر استخر منبع
گنجانده شده برای تنوع مولّد و زبان ترکیب می‌کند؛ جدول~\ref{tab:datastats}
اندازه‌ی پیکره‌ها را خلاصه می‌کند. این پروتکلی سخت‌تر از یک آزمایش دوهدفه‌ی
پیش‌تر این روش است: MLAAD در هر استخر منبع حاضر است (نه فقط به‌عنوان
تشخیصی)، و پیکربندی بهینه‌ساز اصلاح شد (دسته ۳۲، نرخ‌یادگیری
$2\times10^{-4}$؛ پیکربندی پیشین دسته-۶۴/نرخ‌یادگیری-$10^{-4}$ در
جاروب‌های بعدی کم‌آموزش‌دهنده‌ی مدل منبع تشخیص داده شد). اعداد این مقاله
نباید مستقیماً با اعداد آن آزمایش دوهدفه‌ی پیشین مقایسه شوند.

\begin{table}[t]
\centering
\caption{ترکیب پیکره‌ی استخر منبع (تعداد کلیپ پس از تراز کلاس؛ زیرمجموعه‌ی
آب‌پرکرده که واقعاً برای آموزش ذخیره شده، نه اندازه‌ی خام پیکره).}
\label{tab:datastats}
\begin{tabular}{@{}lrr@{}}
\toprule
پیکره & جعلی & واقعی \\
\midrule
ASVspoof2019-LA & ۶۰۰۰ & ۲۵۸۰ \\
LibriSpeech-TTS (D2) & ۲۱۷۳ & ۲۲۷۴ \\
In-the-Wild & ۶۰۰۰ & ۶۰۰۰ \\
عربی (ArAD) & ۶۰۰۰ & ۲۵۷۰ \\
MLAAD (۳۸ زبان، فقط جعلی) & ۸۰۰۰ & --- \\
\bottomrule
\end{tabular}
\end{table}

\textbf{صدای هدفِ نگه‌داشته‌شده هرگز در هیچ استخر منبعی نیست}، از جمله محتوای
هم‌پوشان با زبان MLAAD آن در صورت وجود. درون MLAAD، شش زبان (مالتی،
اسلوونیایی، مجاری، کرواتی، فنلاندی، لیتوانیایی) به‌طور کامل از \emph{هر}
استخر منبع نگه‌داشته می‌شوند، تا بخش~\ref{sec:perlang} تعمیم واقعی زبان
ندیده را بسنجد، نه فقط تعمیم پیکره‌ی ندیده را.

\textbf{نقاط مقایسه.} \emph{فقط-منبع}: آشکارساز فاین‌تیون‌شده، بدون تطبیق.
\emph{Tent}~\cite{tent}: کمینه‌سازی آنتروپی روی همان مجموعه‌ی پارامتر
آموزش‌پذیر روش ما. \emph{فقط-BN}: بازمحاسبه‌ی آماره‌های نرمال‌سازی روی داده‌ی
هدف بدون هیچ گام گرادیانی (حذفی منطبق با بازکالیبراسیون سبک-AdaBN~\cite{adabn}
برای یک معماری LayerNorm). \emph{فقط-خودآموزی}: روش ما با $\lambda{=}0$
(بدون جمله‌ی سازگاری). \emph{روش ما}: روش کامل (بخش III-D). \emph{اوراکل}:
همان معماری با فاین‌تیون روی برچسب‌های هدف برای ۴ دوره---یک کران بالای
نظارت‌شده، نه یک پایه‌ی قابل‌استقرار. \emph{RawNet2Lite}
(بخش~\ref{sec:backbone}): یک ستون‌فقرات پیچشی-بازگشتی از صفر در خانواده‌ی
RawNet2~\cite{rawnet2}، بدون پیش‌آموزش SSL، آموزش‌دیده روی همان استخر منبع
برای همان تعداد دوره، برای آزمودن اینکه آیا استواری میان‌پیکره‌ای واقعاً از
پیش‌سر SSL می‌آید یا از چیزی خاص معماری تجمیع توجهی ما. \emph{DANN}
(بخش~\ref{sec:dann}): آموزش خصمانه‌ی دامنه با بازگشت گرادیان~\cite{dann}
که ورودی‌های \emph{بدون‌برچسب} هدف نگه‌داشته‌شده را به‌عنوان دامنه‌ی خصمانه
در حین آموزش منبع استفاده می‌کند (یعنی همان بودجه‌ی اطلاعاتیِ زمان‌آزمون
روش ما را دارد، اما آن را در زمان آموزش با ۸ دوره بهینه‌سازی خصمانه به‌جای
زمان آزمون خرج می‌کند).

\textbf{معیار.} نرخ خطای برابر (EER) و ROC-AUC. با نرخ‌های مثبت‌کاذب و
منفی‌کاذب $\mathrm{FPR}(\tau),\mathrm{FNR}(\tau)$،
\begin{equation}
\mathrm{EER}=\mathrm{FPR}(\tau^\star)=\mathrm{FNR}(\tau^\star),\quad
\tau^\star:\ \mathrm{FPR}(\tau^\star)=\mathrm{FNR}(\tau^\star),
\end{equation}
که بدون‌آستانه است و پس نسبت به رانش کالیبراسیون بالا حساس نیست.

\textbf{تنظیم.} TTA تراگذر (تطبیق روی استخر هدف بدون‌برچسب، سپس نمره‌دهی
آن) قرارداد استاندارد برای TTA است و چیزی است که مگر ذکر شود گزارش
می‌کنیم؛ افزون بر آن یک بررسی استقرایی (تطبیق روی نیمی از هدف، آزمون روی
نیمه‌ی مجزا) در بخش~\ref{sec:inductive} گزارش می‌کنیم.

\textbf{پیاده‌سازی.} صدای ۱۶کیلوهرتز در ۴ثانیه ذخیره و در زمان آموزش به‌طور
تصادفی به ۳ثانیه بریده می‌شود (تقویت بهره/نویز حفظ‌کننده‌ی برچسب)؛ کل
ذخیره (۴۱۵۹۷ کلیپ) یک‌بار به یک تانسور fp16 واحد مقیم روی GPU رمزگشایی
می‌شود (۵٫۳ گیگابایت)، که کاملاً از DataLoader اجتناب می‌کند---DataLoaderهای
چندکارگر روی برش MIG GPU ما اتمام حافظه‌ی مشترک ایجاد می‌کردند. autocast
bf16؛ Adam، نرخ‌یادگیری $2\times10^{-4}$؛ فاین‌تیون ۴لایه‌ی بالایی XLS-R؛
مدل منبع ۸ دوره آموزش، دسته ۳۲؛ یک برش تک MIG 1g.35gb GPU (حدود ۳۳
گیگابایت). همه‌ی نتایج میانگین$\,\pm\,$انحراف‌معیار روی سه بذر (۰، ۱، ۲)
هستند. با توجه به $n{=}3$، سازگاری علامت به‌ازای بذر («بهبود در $k$ از ۳
بذر») را همراه میانگین گزارش می‌کنیم، نه یک ادعای معناداری آماری.

\section{نتایج}
\subsection{نتایج اصلی میان‌پیکره‌ای}
\label{sec:main-results}
جدول~\ref{tab:main} نتیجه‌ی اصلی است: EER٪/AUC یک-پیکره-بیرون، میانگین روی
سه بذر، در همه‌ی چهار هدف و شش روش. سه الگو برجسته است. نخست، \textbf{روش
ما در سه از چهار هدف برنده یا هم‌تراز است}: ASVspoof2019
($4.76\!\to\!3.69$، هم‌تراز با $3.70$ی Tent در محدوده‌ی نویز، اما با AUC
بالاتر و واریانس بذر‌به‌بذر بسیار کمتر)، In-the-Wild ($12.57\!\to\!10.83$،
کاهش نسبی ۱۳٫۸٪)، و عربی ($22.59\!\to\!20.74$، کاهش نسبی ۸٫۲٪). دوم،
\textbf{روی LibriSpeech-TTS روش خنثی است} ($34.36\!\to\!34.38$، یعنی بدون
بهبود، در محدوده‌ی یک انحراف‌معیار از فقط-منبع، و در $1/3$ بذرها بدتر)؛ به
این در بخش~\ref{sec:gaindecomp} بازمی‌گردیم. سوم، \textbf{فقط-BN روی هر
هدف عددی هم‌ارز فقط-منبع است}، همان‌طور که استدلال LayerNorm در بخش
III-B پیش‌بینی می‌کرد: هیچ آماره‌ی دسته‌ای برای این معماری برای بازمحاسبه
نیست، پس این حذفی یک کنترل منفی تمیز است، نه یک روش رقیب. \textbf{Tent}
بسیار ناپایدار است: روی ASVspoof2019 (جابه‌جایی کوچک) رقابتی است اما روی
عربی به سمت شانس فرومی‌پاشد ($47.66\%$ EER، AUC ۰٫۵۲) و روی In-the-Wild
به‌شدت متغیر است ($26.53\!\pm\!11.46\%$، یعنی گاهی کار می‌کند و گاهی روی
بذرهای \emph{همان} هدف فرومی‌پاشد)---که تأکید می‌کند طراحی برچسب-شبه-مطمئن،
نه TTA به‌طور کلی، است که تطبیق را اینجا پایدار می‌کند. سطر \textbf{اوراکل}
فضای زیاد باقی‌مانده روی هر هدف نشان می‌دهد (مثلاً عربی $20.74\!\to\!9.14\%$
با برچسب‌های هدف)، که تأیید می‌کند هیچ‌کدام از روش‌های بدون‌نظارت به
اشباع آنچه معماری روی این پیکره‌ها می‌تواند برسد نزدیک نیستند.

\begin{table*}[t]
\centering
\caption{EER٪/AUC میان‌پیکره‌ای، میانگین روی ۳ بذر (یک-پیکره-بیرون،
تراگذر؛ EER کمتر بهتر است). بهترین از \{منبع، Tent، ST-only، روش ما\} در
هر ستون \textbf{پررنگ}؛ فقط-BN یک کنترل بی‌اثر است (هم‌ارز منبع، بخش
III-B) و اوراکل یک کران بالای نظارت‌شده است، هر دو از پررنگ‌سازی مستثنا.}
\label{tab:main}
\begin{tabular}{@{}lcccc@{}}
\toprule
روش & ASVspoof2019 & LibriSpeech-TTS (D2) & In-the-Wild & عربی (ArAD) \\
\midrule
فقط-منبع & ۴٫۷۶ / .۹۹۰ & ۳۴٫۳۶ / .۷۲۰ & ۱۲٫۵۷ / .۹۴۴ & ۲۲٫۵۹ / .۸۵۴ \\
بازکالیبراسیون BN~\cite{adabn} & ۴٫۷۶ / .۹۹۰ (بی‌اثر) & ۳۴٫۳۶ / .۷۲۰ (بی‌اثر) & ۱۲٫۵۷ / .۹۴۴ (بی‌اثر) & ۲۲٫۵۹ / .۸۵۴ (بی‌اثر) \\
Tent~\cite{tent} & ۳٫۷۰ / .۹۸۷ & ۳۴٫۷۴ / .۶۹۰ & ۲۶٫۵۳$\pm$۱۱٫۴۶ / .۷۴۳ & ۴۷٫۶۶ / .۵۲۴ \\
فقط-خودآموزی & ۴٫۰۷ / .۹۸۷ & \textbf{۳۳٫۷۹} / .۷۰۸ & ۱۱٫۶۳ / .۹۴۷ & ۲۲٫۰۵ / .۸۵۵ \\
\rowcolor{green!12}\textbf{روش ما (TTA)} & \textbf{۳٫۶۹} / .۹۹۴ & ۳۴٫۳۸ / .۷۲۲ & \textbf{۱۰٫۸۳} / .۹۶۰ & \textbf{۲۰٫۷۴} / .۸۷۳ \\
\midrule
اوراکل (نظارت‌شده روی هدف) & ۰٫۷۰ / ۱٫۰۰۰ & ۱۲٫۵۴ / .۹۵۰ & ۳٫۶۵ / .۹۹۴ & ۹٫۱۴ / .۹۷۰ \\
\bottomrule
\end{tabular}
\end{table*}

\subsection{تجزیه‌ی بهبود: خودآموزی در برابر سازگاری}
\label{sec:gaindecomp}
جدول~\ref{tab:gaindecomp} تغییر کلی EER را به یک گام خودآموزی
(منبع$\to$فقط-خودآموزی) و یک گام سازگاری (فقط-خودآموزی$\to$روش ما)، به‌ازای
هدف تجزیه می‌کند. روی سه هدف، دو مؤلفه مکمل‌اند: خودآموزی به‌تنهایی بیشتر یا
بخشی از بهبود را می‌گیرد، و سازگاری کاهش بیشتری می‌افزاید (بیشترین روی
عربی، بزرگ‌ترین شکاف دامنه: $-1.31$ امتیاز فقط از سازگاری). روی
\textbf{LibriSpeech-TTS، الگو معکوس می‌شود}: خودآموزی به‌تنهایی \emph{کمک
می‌کند} ($-0.57$)، اما افزودن سازگاری $+0.59$ \emph{هزینه} دارد و بهبود را
پاک می‌کند. ما این را همراه با روایت رتبه‌بندی-کالیبراسیون بخش III-C تفسیر
می‌کنیم: LibriSpeech-TTS هدفی است با کمترین AUC منبع ($0.720$، در برابر
$\geq\!0.85$ روی هر هدف دیگر)، یعنی برچسب‌های شبه دنباله‌ی مطمئن مورد
استفاده‌ی هر دو مؤلفه اینجا خودشان کمتر قابل‌اعتمادند؛ جمله‌ی سازگاری، که
پیش‌بینی‌ها روی نماهای تقویت‌شده و تمیز را به‌هم می‌کشد صرف‌نظر از اینکه
کدام درست است، به‌نظر می‌رسد به‌جای اصلاح این عدم‌قابلیت‌اعتماد، آن را
تشدید می‌کند وقتی سیگنال رتبه‌بندی زیربنایی این‌قدر ضعیف است. این محدودیتی
واقعی از روش با پیکربندی فعلی است، نه مصنوع تنظیمی که انتخاب نکردیم اصلاح
کنیم: جاروب فراپارامتر بخش~\ref{sec:sweep} (روی دو هدف اصلی اجرا شده) از
پیش نشان داده بود چندک بهینه‌ی دنباله‌ی مطمئن وابسته به هدف است، و این
نتیجه آن یافته را به خود وزن سازگاری هم گسترش می‌دهد.

\begin{table}[t]
\centering
\caption{تجزیه‌ی بهبود، تغییر میانگین EER٪ روی ۳ بذر. $\Delta_1$:
منبع$\to$فقط-خودآموزی. $\Delta_2$: فقط-خودآموزی$\to$روش ما (سهم حاشیه‌ای
جمله‌ی سازگاری).}
\label{tab:gaindecomp}
\begin{tabular}{@{}lccc@{}}
\toprule
هدف & EER منبع & $\Delta_1$ (خودآموزی) & $\Delta_2$ (سازگاری) \\
\midrule
ASVspoof2019 & ۴٫۷۶ & $-0.69$ & $-0.39$ \\
LibriSpeech-TTS & ۳۴٫۳۶ & $-0.57$ & \cellcolor{red!10}$+0.59$ \\
In-the-Wild & ۱۲٫۵۷ & $-0.94$ & $-0.79$ \\
عربی & ۲۲٫۵۹ & $-0.54$ & $-1.31$ \\
\bottomrule
\end{tabular}
\end{table}

\subsection{آیا ستون‌فقرات اهمیت دارد؟ مقایسه‌ی غیر-SSL}
\label{sec:backbone}
جدول~\ref{tab:backbone} مدل منبع XLS-R فاین‌تیون‌شده را در برابر
RawNet2Lite، یک ستون‌فقرات از صفر بدون پیش‌آموزش خودنظارتی، آموزش‌دیده به‌طور
یکسان روی همان استخر منبع مقایسه می‌کند. شکاف چشمگیر است: RawNet2Lite هرگز
از AUC ۰٫۶۵ روی هیچ هدفی فراتر نمی‌رود و روی سه از چهار هدف در یا زیر شانس
است (AUC عربی ۰٫۴۱، یعنی بیشتر اوقات \emph{وارون}با برچسب واقعی همبسته
است، نشانه‌ی بیش‌برازش شدید به مصنوعات پیکره‌ی منبع به‌جای ویژگی‌های
مرتبط با جعل). این تأیید می‌کند که کارآمدی میان‌پیکره‌ای در تنظیمات ما از
پیش‌آموزش SSL آغاز می‌شود، نه از گام TTA: TTA می‌تواند چند امتیاز EER را از
مدل منبعی که از پیش تا حدی تعمیم می‌یابد (AUC $\gtrsim\!0.72$) بازیابد،
اما جایگزینی برای نمایشی که ابتدا منتقل می‌شود نیست، و ستون‌فقراتی به این
فاصله از شانس را نجات نمی‌دهد.

\begin{table}[t]
\centering
\caption{EER٪/AUC فقط-منبع بر حسب ستون‌فقرات، میانگین روی ۳ بذر (تراگذر،
بدون تطبیق در هیچ سطر).}
\label{tab:backbone}
\begin{tabular}{@{}lcc@{}}
\toprule
هدف & XLS-R (ما) & RawNet2Lite~\cite{rawnet2} \\
\midrule
ASVspoof2019 & ۴٫۷۶ / .۹۹۰ & ۴۷٫۴۷ / .۵۳۹ \\
LibriSpeech-TTS & ۳۴٫۳۶ / .۷۲۰ & ۵۱٫۱۴ / .۴۸۳ \\
In-the-Wild & ۱۲٫۵۷ / .۹۴۴ & ۳۹٫۳۵ / .۶۴۹ \\
عربی & ۲۲٫۵۹ / .۸۵۴ & ۵۷٫۰۱ / .۴۱۱ \\
\bottomrule
\end{tabular}
\end{table}

\subsection{مقایسه با آموزش خصمانه‌ی دامنه در زمان آموزش}
\label{sec:dann}
جدول~\ref{tab:dann} روش زمان‌آزمون ما را در برابر DANN~\cite{dann}، که تا
همگرایی آموزش دیده (۸ دوره بهینه‌سازی خصمانه) و با دسترسی به ورودی‌های
بدون‌برچسب هدف نگه‌داشته‌شده \emph{در حین آموزش منبع}---دسترسی به‌مراتب
ممتازتر به هدف نسبت به روش ما که هدف را فقط در زمان آزمون، پس از پایان
آموزش منبع، می‌بیند---مقایسه می‌کند. با این حال، DANN در سه از چهار هدف
عملکردی ضعیف‌تر از روش ما دارد (روی ASVspoof2019 نزدیک است، جایی که هر دو
از پیش نزدیک اوراکل‌اند) و روی عربی به‌طور محسوسی بدتر است ($29.51\%$ در
برابر $20.74\%$ ما). همچنین تقریباً ۳ برابر گران‌تر است: حدود ۳۰ دقیقه آموزش
خصمانه به‌ازای هر بذر-هدف در برابر حدود ۹ دقیقه تطبیق زمان‌آزمون. دو شیوه‌ی
شکست در گزارش آموزش میان اهداف تکرار می‌شود: وزن تلفات خصمانه
$\lambda_{\mathrm{dann}}$ زمان‌بندی خود را ظرف ۳--۴ دوره اشباع می‌کند، و
پس از آن سهم گرادیانی طبقه‌بند دامنه عملاً برای باقی آموزش ثابت است، و---سازگار
با آزمایش اولیه و کوچک‌تر ما با DA-RAD که همان سازوکار بازگشت گرادیان را
استفاده می‌کرد---آموزش خصمانه ناپایدارترین روش از نظر بذر است که می‌آزماییم
(EER ASVspoof2019 میان بذرها از $3.23$ تا $10.43\%$ است، گستره‌ای ۳ برابر).
این خط شواهد دوم و مستقل (فراتر از حذفی DA-RAD پیشین) می‌افزاید که اهداف
خصمانه‌ی زمان‌آموزش دامنه مسیری قابل‌اعتماد برای استواری میان‌پیکره‌ای در
این وظیفه نیستند، و اینکه سازوکار سبک‌تر زمان‌آزمون نه‌تنها ارزان‌تر بلکه
مؤثرتر هم هست.

\begin{table}[t]
\centering
\caption{تطبیق زمان‌آزمون در برابر آموزش خصمانه‌ی دامنه در زمان آموزش
(DANN)، میانگین EER٪/AUC روی ۳ بذر و میانگین دقیقه‌ی واقعی به‌ازای
بذر-هدف.}
\label{tab:dann}
\begin{tabular}{@{}lccc@{}}
\toprule
هدف & منبع & DANN~\cite{dann} (۳۰ دقیقه) & روش ما (۹ دقیقه) \\
\midrule
ASVspoof2019 & ۴٫۷۶ / .۹۹۰ & ۶٫۶۱ / .۹۸۰ & \textbf{۳٫۶۹} / .۹۹۴ \\
LibriSpeech-TTS & ۳۴٫۳۶ / .۷۲۰ & ۳۴٫۸۱ / .۷۲۴ & \textbf{۳۴٫۳۸} / .۷۲۲ \\
In-the-Wild & ۱۲٫۵۷ / .۹۴۴ & ۱۱٫۹۵ / .۹۳۱ & \textbf{۱۰٫۸۳} / .۹۶۰ \\
عربی & ۲۲٫۵۹ / .۸۵۴ & ۲۹٫۵۱ / .۷۲۴ & \textbf{۲۰٫۷۴} / .۸۷۳ \\
\bottomrule
\end{tabular}
\end{table}

\subsection{حساسیت فراپارامتری}
\label{sec:sweep}
سه دستگیره‌ی تطبیق---چندک دنباله‌ی مطمئن $q$، دوره‌های تطبیق $E$، و وزن
سازگاری $\lambda$---را یکی‌یکی حول پیش‌فرض‌ها ($q{=}0.3$، $E{=}4$،
$\lambda{=}0.3$) که در جای دیگر مقاله استفاده می‌شود جاروب می‌کنیم. این
جاروب پیش از گسترش چهارهدفه بود و روی بذر ۰ فقط دو هدف اصلی (In-the-Wild،
عربی) اجرا شد، به دلایل بودجه‌ی محاسباتی؛ آن را همان‌طور که بود گزارش
می‌کنیم نه با بازاجرا، چون نتیجه‌گیری‌های کیفی (زیر) با نتیجه‌ی جدیدتر
بخش~\ref{sec:gaindecomp} روی هدف سوم و مستقلی تأیید می‌شوند
(جدول~\ref{tab:sweep}).\footnote{یک جاروب پیشین یک جریان RNG پیوسته را
میان پیکربندی‌ها مشترک می‌کرد، پس نقاط مجاور تا حدی با نویز ترتیب دسته
تفاوت داشتند نه فقط با دستگیره؛ هر نقطه‌ی گزارش‌شده اینجا بلافاصله پیش از
تطبیق \texttt{torch.manual\_seed(0)} را دوباره بذر می‌زند، و اثر دستگیره
را جدا می‌کند.} دو یافته‌ی سازگار. نخست، \emph{دوره‌های تطبیق بیشتر روی هر
دو هدف کمک می‌کند} ($E{=}8$ از $E{=}4$ به میزان ۱٫۲ امتیاز روی In-the-Wild
و ۱٫۶ روی عربی بهتر است): پیش‌فرض ۴دوره‌ای ما انتخابی محافظه‌کارانه و
منطبق‌با-محاسبات برای جدول اصلی بود، نه بهینه‌ی تنظیم‌شده، و استقراری که
بتواند گام‌های تطبیق بیشتری بپردازد باید انتظار بهبود بیشتری داشته باشد.
دوم، \emph{وزن سازگاری $\lambda{>}0$ روی این دو هدف یکسان لازم است}: تنظیم
$\lambda{=}0$ بدترین پیکربندی جاروب روی هر دو است. بخش~\ref{sec:gaindecomp}
نشان می‌دهد این یافته‌ی دوم به LibriSpeech-TTS \emph{تعمیم نمی‌یابد}،
جایی که $\lambda{=}0$ در واقع بهتر است---پس «سازگاری همیشه کمک می‌کند»
باید وابسته‌به‌هدف خوانده شود، نه همگانی، برخلاف چیزی که جاروب دوهدفه‌ی
اصلی به‌تنهایی پیشنهاد می‌داد. چندک بهینه‌ی دنباله‌ی مطمئن هم وابسته‌به‌هدف
است ($q{=}0.3$ روی In-the-Wild در برابر $q{=}0.2$ روی عربی، هدف با
جابه‌جایی بزرگ‌تر)، که نشان می‌دهد کسر پیش‌بینی‌های قابل‌اعتماد دنباله با
رشد شکاف دامنه تا حدی کوچک می‌شود. ما یک پیکربندی ثابت را برای کل مقاله
حفظ می‌کنیم تا مقایسه‌ی منصفانه‌ی میان‌هدفی داشته باشیم به‌جای تنظیم
به‌ازای هدف؛ تنظیم به‌ازای هدف جهتی طبیعی برای استقراری است که بتواند یک
سیگنال اعتبارسنجی داشته باشد.

\begin{figure}[t]
\centering
\safefig{0.86\columnwidth}{fig_auc_gain.png}
\caption{بهبود EER (منبع$-$روش ما) در برابر AUC منبع پیش از تطبیق، یک نقطه
به‌ازای (هدف، بذر)، هر ۱۲ اجرای یک-پیکره-بیرون. روند مثبت
($r{=}{+}0.41$) از روایت پیش‌شرط رتبه‌بندی بخش III-C پشتیبانی می‌کند:
LibriSpeech-TTS (قرمز)، هدف با کمترین AUC، در یا زیر بهبود صفر خوشه
می‌بندد، در حالی که هدف‌های AUC بالاتر بالای آن خوشه می‌بندند---هرچند این
رابطه در $n{=}12$ نویزی است و یک بذر LibriSpeech-TTS بهبود منفی کوچکی حتی
نسبت به AUC خودش نشان می‌دهد.}
\label{fig:aucgain}
\end{figure}

\begin{table}[t]
\centering
\caption{جاروب فراپارامتر (بذر ۰، زیرمجموعه‌ی دوهدفه، EER٪/AUC؛ سطر
پیش‌فرض $q{=}0.3,E{=}4,\lambda{=}0.3$ به رنگ \colorbox{green!12}{سبز}).}
\label{tab:sweep}
\begin{tabular}{@{}lcc@{}}
\toprule
تنظیم & In-the-Wild & عربی \\
\midrule
فقط-منبع & ۱۸٫۷۳ / .۹۰۰ & ۲۸٫۵۶ / .۷۸۱ \\
\midrule
$q=0.1$ & ۱۴٫۸۷ / .۹۳۳ & ۲۵٫۱۴ / .۸۲۵ \\
$q=0.2$ & ۱۴٫۷۰ / .۹۳۳ & \textbf{۲۳٫۰۵} / \textbf{.۸۴۹} \\
\rowcolor{green!12}$q=0.3$ (پیش‌فرض) & \textbf{۱۴٫۲۵} / \textbf{.۹۳۴} & ۲۴٫۷۴ / .۸۳۶ \\
$q=0.4$ & ۱۴٫۴۷ / .۹۳۹ & ۲۴٫۵۱ / .۸۳۰ \\
\midrule
$E=2$ & ۱۵٫۳۸ / .۹۲۳ & ۲۶٫۹۹ / .۸۰۹ \\
\rowcolor{green!12}$E=4$ (پیش‌فرض) & ۱۴٫۲۵ / .۹۳۴ & ۲۴٫۷۴ / .۸۳۶ \\
$E=8$ & \textbf{۱۳٫۰۳} / \textbf{.۹۴۵} & \textbf{۲۳٫۱۲} / \textbf{.۸۵۲} \\
\midrule
$\lambda=0$ (بدون سازگاری) & ۱۵٫۷۵ / .۹۱۹ & ۲۵٫۴۶ / .۸۲۲ \\
\rowcolor{green!12}$\lambda=0.3$ (پیش‌فرض) & ۱۴٫۲۵ / .۹۳۴ & \textbf{۲۴٫۷۴} / \textbf{.۸۳۶} \\
$\lambda=1.0$ & \textbf{۱۳٫۹۸} / \textbf{.۹۳۷} & ۲۵٫۲۸ / .۸۳۲ \\
\bottomrule
\end{tabular}
\end{table}

\subsection{بررسی استقرایی}
\label{sec:inductive}
برای رد به‌خاطرسپاری تراگذر، روی نیمی از هر هدف تطبیق می‌دهیم و روی نیمه‌ی
مجزا ارزیابی می‌کنیم (جدول~\ref{tab:inductive}). بهبود روی سه هدفی که
نتیجه‌ی تراگذر مثبت است پابرجا می‌ماند: EER روی ASVspoof2019
($4.93\!\to\!4.16$، $3/3$ بذر)، In-the-Wild ($12.66\!\to\!11.13$، $3/3$
بذر)، و عربی ($22.29\!\to\!21.14$، $3/3$ بذر) بهبود می‌یابد---در هر مورد در
محدوده‌ی یک امتیاز از نتیجه‌ی تراگذر، پس بهبود ساختار سطح‌دامنه‌ی آموخته‌شده
در حین تطبیق را بازتاب می‌دهد، نه به‌خاطرسپاری کلیپ‌های خاص نمره‌داده‌شده.
روی LibriSpeech-TTS نتیجه‌ی استقرایی یافته‌ی خنثای تراگذر را بازتاب می‌دهد
(منبع $34.49\!\to\!$ روش ما $34.70$، یعنی کمی بدتر، سازگار با یافته‌ی
$\Delta_2{>}0$ بخش~\ref{sec:gaindecomp})---خنثی بودن روش روی این هدف مصنوع
ارزیابی تراگذر هم نیست.

\begin{table}[t]
\centering
\caption{بررسی استقرایی: تطبیق روی نیمی از هدف، ارزیابی روی نیمه‌ی مجزا.
میانگین EER٪ روی ۳ بذر.}
\label{tab:inductive}
\begin{tabular}{@{}lcc@{}}
\toprule
هدف & منبع (نیمی) & روش ما (نیمی، ارزیابی مجزا) \\
\midrule
ASVspoof2019 & ۴٫۹۳ & \textbf{۴٫۱۶} \\
LibriSpeech-TTS & ۳۴٫۴۹ & ۳۴٫۷۰ \\
In-the-Wild & ۱۲٫۶۶ & \textbf{۱۱٫۱۳} \\
عربی & ۲۲٫۲۹ & \textbf{۲۱٫۱۴} \\
\bottomrule
\end{tabular}
\end{table}

\subsection{آیا تطبیق واگرایی دامنه را کاهش می‌دهد؟}
\label{sec:divergence}
یک فرضیه‌ی طبیعی این است که تطبیق به همان شیوه‌ای کار می‌کند که روش‌های DG
خصمانه صریحاً برای آن طراحی شده‌اند: با تمایزناپذیر کردن منبع و هدف در فضای
نهفته. جدول~\ref{tab:divergence} این را مستقیماً با فاصله‌ی $\mathcal{A}$ی
جانشین بخش~\ref{sec:method-divergence} و یک کاوشگر منشأ دامنه‌ی خطی منطبق،
پیش و پس از تطبیق روی هر هدف، می‌آزماید. روی \emph{هر} هدف، هر دو معیار در
جهت \emph{مخالف} چیزی که فرضیه‌ی تغییرناپذیری دامنه پیش‌بینی می‌کند حرکت
می‌کنند: $\hat d_{\mathcal{A}}$ پس از تطبیق کمی افزایش می‌یابد (مثلاً عربی
$1.545\!\to\!1.606$)، و دقت کاوشگر دامنه در تشخیص نهفتگی‌های منبع از هدف هم
کمی بالا می‌رود (مثلاً عربی $87.9\%\!\to\!89.1\%$). این حتی روی عربی و
ASVspoof2019، جایی که EER به‌طور چشمگیری بهبود می‌یابد، برقرار است---پس
بهبود EER و افزایش واگرایی دامنه \emph{همزمان} رخ می‌دهند. نتیجه می‌گیریم
که روش ما با پاک کردن هویت دامنه از نمایش کار \emph{نمی‌کند}؛ با
بازشکل‌دهی مرز تصمیم با استفاده از برچسب‌های شبه رتبه-قابل‌اعتماد کار
می‌کند، سازگار با روایت رتبه‌بندی-کالیبراسیون بخش III-C، نه با روایت
هم‌ترازی نمایش. این نتیجه‌ی منفی مفیدی برای ادبیات DG اعمال‌شده به این
وظیفه است: نشان می‌دهد جریمه‌ی صریح جداپذیری دامنه (کاری که DANN می‌کند)
کمیتی را هدف می‌گیرد که واقعاً گلوگاه نیست، شاهدی مستقل و همسو با چرایی
عملکرد ضعیف‌تر DANN در بخش~\ref{sec:dann} با وجود بهینه‌سازی دقیقاً همین
کمیت.

\begin{table}[t]
\centering
\caption{فاصله‌ی $\mathcal{A}$ی جانشین و دقت کاوشگر منشأ دامنه‌ی خطی، پیش
در برابر پس از تطبیق (بذر ۰). هر دو روی هر هدف پس از تطبیق کمی افزایش
می‌یابند---جداپذیری دامنه کاهش نمی‌یابد.}
\label{tab:divergence}
\begin{tabular}{@{}lcccc@{}}
\toprule
& \multicolumn{2}{c}{$\hat d_{\mathcal{A}}$} & \multicolumn{2}{c}{دقت کاوشگر دامنه (٪)} \\
هدف & منبع & تطبیق‌یافته & منبع & تطبیق‌یافته \\
\midrule
ASVspoof2019 & ۱٫۵۷۲ & ۱٫۶۰۹ & ۸۶٫۶ & ۸۹٫۹ \\
LibriSpeech-TTS & ۱٫۴۶۰ & ۱٫۵۰۴ & ۸۶٫۴ & ۸۶٫۶ \\
In-the-Wild & ۱٫۳۳۵ & ۱٫۳۵۷ & ۸۰٫۴ & ۸۱٫۶ \\
عربی & ۱٫۵۴۵ & ۱٫۶۰۶ & ۸۷٫۹ & ۸۹٫۱ \\
\bottomrule
\end{tabular}
\end{table}

\subsection{هندسه‌ی نهفتگی: کلاس فضا را سازمان می‌دهد، نه دامنه}
شکل~\ref{fig:embed} ۷۵۰ نهفتگی منبع و ۷۵۰ نهفتگی هدف را با t-SNE برای هر
چهار هدف، پیش و پس از تطبیق، رنگ‌کدگذاری‌شده بر اساس منشأ منبع-در-برابر-هدف
تصویرسازی می‌کند. در هر پانل، نقاط منبع و هدف به‌طور کامل \emph{درهم‌آمیخته}‌اند
سراسر نمودار به‌جای تشکیل دو جزیره‌ی دامنه‌ی جدا---دو کمان بزرگ قابل‌مشاهده
در هر پانل متناظر با ساختار کلاس واقعی/جعلی است (با ارجاع‌متقابل به
برچسب‌های کلاس تأیید شد)، و منبع/هدف \emph{درون} هر کمان درهم‌می‌آمیزند.
این همتای کیفی جدول~\ref{tab:divergence} است: نمایش از پیش عمدتاً بر اساس
\emph{کلاس} سازمان یافته، نه پیکره‌ی منشأ، هم پیش و هم پس از تطبیق، که
سازگار است با چرایی بقای رتبه‌بندی میان‌پیکره‌ای زیر جابه‌جایی (بخش III-C)
و با چرایی نیاز نداشتن تطبیق به (و ندادن) حذف اطلاعات دامنه برای کار کردن.

\begin{figure}[t]
\centering
\safefig{\columnwidth}{fig_tsne_multitarget.png}
\caption{t-SNE از ۷۵۰ نهفتگی منبع + ۷۵۰ نهفتگی هدف برای هر چهار هدف
(ستون‌ها)، پیش (ردیف بالا) در برابر پس از تطبیق (ردیف پایین)؛
آبی$=$منبع، نارنجی$=$هدف. در هر پانل، منبع و هدف درون هریک از دو کمان
ساختاریافته‌ی کلاس درهم‌می‌آمیزند به‌جای جدا شدن به خوشه‌های دامنه‌ویژه، هم
پیش و هم پس از تطبیق---منطبق با نتیجه‌ی فاصله‌ی $\mathcal{A}$ی جانشین که
جداپذیری دامنه با تطبیق کاهش نمی‌یابد (جدول~\ref{tab:divergence}).}
\label{fig:embed}
\end{figure}

\subsection{تعمیم چندزبانه: بازخوانی به‌ازای زبان}
\label{sec:perlang}
چون MLAAD کلیپ‌های جعلی را در ۳۸ زبان به هر استخر منبع می‌افزاید
(بخش IV)، می‌توانیم مستقیماً بیازماییم که آیا بازخوانی کلیپ جعلی
آشکارساز به حضور یا عدم‌حضور یک زبان در آموزش منبع بستگی دارد. شش زبان
(مالتی، اسلوونیایی، مجاری، کرواتی، فنلاندی، لیتوانیایی) از هر استخر منبع
نگه‌داشته می‌شوند؛ ۳۲ زبان باقی‌مانده در دست‌کم یک استخر منبع حاضرند.
میانگین‌گیری روی هر چهار مدل منبع و هر دو گروه (شکل~\ref{fig:perlang})،
میانگین بازخوانی کلیپ جعلی ۹۳٫۹٪ روی زبان‌های دیده و ۹۳٫۳٪ روی زبان‌های
ندیده است---شکافی زیر یک امتیاز، کاملاً در محدوده‌ی پراکندگی به‌ازای زبان
($\sigma\!\approx\!7$ امتیاز) مشاهده‌شده درون هر گروه. تنها ناهنجار،
مالتی روی مدل منبع In-the-Wild است (بازخوانی ۶۳٫۰٪)؛ هیچ زبان دیگری زیر
۶۷٪ نمی‌افتد. این نتیجه‌ای دلگرم‌کننده برای استقرار است: استواری چندزبانه‌ی
آشکارساز صرفاً مصنوع از پیش دیدن مصنوعات ترکیب‌سازی یک زبان خاص در حین
آموزش نیست، که شیوه‌ی شکست نگران‌کننده‌تر برای سامانه‌ای است که قرار است
به زبان‌های واقعاً تازه در زمان استقرار تعمیم یابد.

\begin{figure}[t]
\centering
\safefig{\columnwidth}{fig_perlang.png}
\caption{بازخوانی کلیپ جعلی MLAAD بر حسب زبان، میانگین‌گیری‌شده روی چهار
مدل منبع، مرتب‌شده صعودی. میله‌های نارنجی شش زبان نگه‌داشته‌شده از هر
استخر منبع‌اند؛ میله‌های سبز در دست‌کم یک استخر منبع حاضر بودند. خط
چین‌دار میانگین کلی است ($93.8\%$). زبان‌های ندیده نظام‌مند بدتر از دیده‌ها
نیستند، به‌جز یک ناهنجار تک (مالتی).}
\label{fig:perlang}
\end{figure}

\subsection{چه چیزی کار نکرد}
آموزش خصمانه‌ی دامنه در زمان آموزش (بازگشت گرادیان روی هویت پیکره) به‌همراه
یک تلفات تباینی واقعی-لنگرشده---رویکرد اولیه‌ی ما، ارزیابی‌شده روی یک
رمزگذار منجمد و سپس فاین‌تیون‌شده در آزمایش‌های پیشین---EER میان‌پیکره‌ای را
بهبود نداد و ناپایدار بود (آموزش خصمانه در پیکربندی رمزگذار-فاین‌تیون‌شده
به شانس فروپاشید). پایه‌ی DANN بخش~\ref{sec:dann}، که تا همگرایی به‌جای
فقط آزمایشی آموزش دید، این را با یک پیاده‌سازی دوم و مستقل تأیید می‌کند: با
وجود ۳ برابر محاسبات و دسترسی ممتاز به ورودی‌های هدف در حین آموزش، در سه از
چهار هدف عملکردی ضعیف‌تر از روش ما دارد. تحلیل واگرایی دامنه‌ی
بخش~\ref{sec:divergence} تبیینی پسینی برای هر دو نتیجه ارائه می‌دهد: نمایش
از ابتدا عمدتاً بر اساس دامنه سازمان نیافته بود (بخش‌های ۵٫۸--۵٫۹)، پس حذف
صریح و خصمانه‌ی هویت دامنه سیگنال مفید کمی برای عمل داشت و عمدتاً
بهینه‌سازی را ناپایدار کرد، در حالی که جمله‌ی سازگاری ما (معادله‌ی ۷) هیچ
فرض تغییرناپذیری دامنه‌ای نمی‌کند. درون روش خودمان هم، خودآموزی-با-سازگاری
همیشه بهترین پیکربندی نیست: بخش~\ref{sec:gaindecomp} نشان می‌دهد
فقط-خودآموزی روی LibriSpeech-TTS از روش کامل بهتر است، و ما این را با
تنظیم $\lambda$ به‌ازای هدف پس از واقعه پنهان نمی‌کنیم.

\section{بحث و محدودیت‌ها}
روش در برچسب‌ها بدون‌نظارت است اما \emph{تراگذر} است: ورودی‌های هدف
بدون‌برچسب را در زمان تطبیق استفاده می‌کند، تنظیمی استاندارد در TTA که
صراحتاً بیان می‌کنیم؛ بررسی استقرایی (بخش~\ref{sec:inductive}) نگرانی
به‌خاطرسپاری را برطرف می‌کند و روی هر هدف از نزدیک اعداد تراگذر را دنبال
می‌کند، از جمله بازتولید نتیجه‌ی خنثای LibriSpeech-TTS به‌جای پنهان کردن
آن. رویکرد نسبت به جایگزین‌هایی که با آن‌ها مقایسه می‌کنیم ارزان است (تطبیق
یک مدل آموزش‌دیده در چند دقیقه، در برابر حدود ۳۰ دقیقه برای DANN) و نسبت
به انتخاب پیش‌سر SSL متعامد است، هرچند بخش~\ref{sec:backbone} نشان می‌دهد
جایگزینی برای پیش‌آموزش SSL نیست. ما محدودیت‌های روش را صراحتاً برمی‌شماریم
به‌جای پنهان کردن آن‌ها درون روایت نتایج، چون ترجیح می‌دهیم خواننده آن‌ها را
اینجا بیابد تا خودش کشف کند.

\begin{enumerate}
\item \textbf{فراپارامترهای ثابت، نه تنظیم به‌ازای هدف.} هر نتیجه‌ی این
مقاله از یک پیکربندی ($q{=}0.3$، $E{=}4$، $\lambda{=}0.3$) در هر چهار هدف
استفاده می‌کند، انتخاب‌شده بدون دسترسی به برچسب‌های هدف، دقیقاً چون یک
استقرار واقعی هیچ سیگنال اعتبارسنجی برای تنظیم در برابر آن ندارد
(بخش~\ref{sec:sweep}). جاروب و تجزیه‌ی بهبود (جدول~\ref{tab:gaindecomp})
نشان می‌دهند این پیکربندی ثابت همگانی-بهینه نیست---$\lambda{=}0$ روی
LibriSpeech-TTS بهتر بود، و $q{=}0.2$ روی عربی از پیش‌فرض $q{=}0.3$ ما
بهتر بود---که دقیقاً همان حساسیتی است که یک عملگر بدون برچسب نمی‌تواند
اصلاح کند. این حساسیت را به‌عنوان ویژگی توصیف‌شده‌ی روش گزارش می‌کنیم، نه
شاهدی بر کم‌تنظیم بودن: جایگزینِ گزارش اعداد اوراکل‌تنظیم‌شده به‌ازای هدف
به‌عنوان نتیجه‌ی اصلی، آنچه بدون برچسب در زمان استقرار دست‌یافتنی است را
بیش از واقع نشان می‌دهد.

ما به‌جای فقط گزارش این مسئله، تلاش کردیم آن را رفع کنیم: یک اکتشافی
بدون‌برچسب که $\lambda$ را فقط از نمرات بدون‌برچسب هدف انتخاب می‌کند، با
استفاده از جدایی دنباله‌ی مطمئن ($\mathrm{mean}(s\mid s\!\geq\!Q_{0.7}) -
\mathrm{mean}(s\mid s\!\leq\!Q_{0.3})$، محاسبه‌پذیر بدون هیچ برچسب هدف) به
عنوان یک جانشین برای اینکه آیا رتبه‌بندی مدل منبع روی آن هدف قابل‌اعتماد
است. با اعتبارسنجی پسینی در برابر $\Delta_2$ی واقعی روی هر ۱۲ نقطه‌ی
(هدف، بذر)، این روش \emph{شکست می‌خورد}، و در جهت اشتباه شکست می‌خورد
($r{=}{+}0.74$ میان جانشین و $\Delta_2$، جایی که یک اکتشافی مفید به
$r{<}0$ نیاز دارد): LibriSpeech-TTS، تنها هدفی که سازگاری زیان‌بار است،
\emph{بیشترین} جدایی دنباله را از هر هدف دارد ($0.92$--$0.99$)، نه کمترین
را. ما این را سازگار با، نه در تناقض با، بخش III-C می‌خوانیم: نمرات خام
آشکارساز ما روی \emph{هر} هدف نزدیک افراط‌ها متمرکز می‌شوند صرف‌نظر از
اینکه آن اطمینان درست باشد یا نه، پس جانشینی که فقط از اطمینان نمره
ساخته شده نمی‌تواند مطمئن‌و‌درست را از مطمئن‌و‌بدکالیبره تشخیص دهد، که
دقیقاً همان شیوه‌ی شکستی است که LibriSpeech-TTS نمایندگی می‌کند. یک سیگنال
بدون‌برچسب کارآمد باید چیزی نزدیک‌تر به درستی کالیبراسیون را اندازه بگیرد
تا به اطمینان، که همچنان باز است.

\item \textbf{روش به مدل منبعی با رتبه‌بندی از پیش قابل‌اعتماد نیاز دارد.}
سازوکار بخش III-C نیاز دارد برچسب‌های شبه دنباله‌ی مطمئن قابل‌اعتماد
باشند، که به‌نوبه‌ی خود نیاز دارد AUC منبع روی هدف پیش از آغاز تطبیق
به‌اندازه‌ی کافی بالا باشد. روی هر $4\times3{=}12$ نقطه‌ی (هدف، بذر)، AUC
منبع و بهبود EER حاصل مثبت همبسته‌اند ($r{=}{+}0.41$،
شکل~\ref{fig:aucgain})؛ تنها هدفی که روش روی آن خنثی است
(LibriSpeech-TTS) همان هدفی است با کمترین AUC منبع ($0.72$، در برابر
$\geq\!0.85$ روی سه هدف دیگر). ما این را شاهدی \emph{به‌سود} سازوکار
پیشنهادی خود می‌خوانیم---روش دقیقاً جایی شکست می‌خورد که پیش‌شرط خودش
ضعیف‌ترین است، نه دلبخواه---اما با تنها چهار هدف این یک رابطه‌ی محتمل
است، نه یک قاعده‌ی پیش‌بین اعتبارسنجی‌شده، و یک نقطه‌ی آغاز از صفر یا
AUC-پایین رژیمی نیست که در حال حاضر بتوانیم نجات دهیم.

\item \textbf{آموزش مدل منبع کاملاً پایدار از نظر بذر نیست.} EER منبع
بذر‌به‌بذر روی In-the-Wild از $9.4\%$ تا $13.2\%$ و روی ASVspoof2019 از
$3.5\%$ تا $6.0\%$ زیر دستور ما گستره دارد---گستره‌هایی وسیع‌تر از آنچه
می‌خواستیم، و یک عامل مخدوش‌کننده‌ی بالقوه هنگام جدا کردن اثر حاشیه‌ای
روش تطبیق از نویز آموزش منبع. دو چیز آسیب این را به ادعاهای ما محدود
می‌کند: رتبه‌بندی کیفی روش‌ها (روش ما~$<$~خودآموزی~$<$~منبع~$\ll$~Tent)
با وجود این واریانس میان بذرها سازگار می‌ماند، و ناپایداری بذر‌به‌بذر خود
Tent (بخش~\ref{sec:main-results}) به‌طور قابل‌توجهی بزرگ‌تر از واریانس
آموزش منبع است. کاهش این واریانس (مثلاً یک زمان‌بندی نرخ‌یادگیری تنظیم‌شده،
یا تعداد بذر بیشتر از آنچه بودجه‌ی محاسباتی ما اجازه داد) کار آینده است.

\item \textbf{توان آماری با مقیاس محدود است.} چهار هدف و سه بذر سازگاری
علامت و یک سازوکار محتمل را برقرار می‌کنند، نه یک آزمون آماری دقیق.
به‌ویژه، نمی‌توانیم تعیین کنیم آیا نتیجه‌ی خنثای LibriSpeech-TTS زیر
ترکیب استخر منبع متفاوت، تفکیک تصادفی متفاوتی از آن پیکره، یا هدف پنجمی با
AUC منبع مشابه پایین بازتولید می‌شود یا نه. ما اعداد به‌ازای هدف و به‌ازای
بذر را در سراسر (جدول‌های~\ref{tab:main}--\ref{tab:divergence}) گزارش
می‌کنیم به‌جای یک عدد اصلی تجمیع‌شده میان اهداف، دقیقاً تا نتیجه‌ی
LibriSpeech-TTS نتواند میانگین‌گیری‌شده و پنهان شود؛ یک بازتولید در مقیاس
بزرگ‌تر با اهداف و بذرهای بیشتر مستقیم‌ترین راه برای رفع این محدودیت است.

\item \textbf{نتیجه به پیش‌آموزش SSL بستگی دارد، نه فقط به TTA.}
بخش~\ref{sec:backbone} نشان می‌دهد یک ستون‌فقرات غیر-SSL از صفر
(RawNet2Lite) آموزش‌دیده به‌طور یکسان روی همان استخر منبع در یا زیر شانس
در سه از چهار هدف است (AUC تا $0.41$ روی عربی). این با محدودیت ۲ بالا
سازگار است---ستون‌فقراتی بدون رتبه‌بندی قابل‌اعتماد چیزی برای برچسب‌گذاری
شبه دنباله‌ی مطمئن برای بهره‌بردن ندارد---اما همچنین یعنی سهم ما باید
\emph{یک لایه‌ی تطبیق سبک روی نمایش SSL از پیش قابل‌انتقال} خوانده شود، نه
یک دستور همه‌کاره برای استواری میان‌پیکره‌ای مستقل از انتخاب ستون‌فقرات.
RawNet2Lite اینجا کوچک است (۱٫۲ میلیون پارامتر) و از صفر فقط برای هشت
دوره آموزش دیده، پس نمی‌توانیم به‌طور کامل «پیش‌آموزش SSL لازم است» را از
«این پایه‌ی خاص غیر-SSL کم‌آموزش دیده بود» جدا کنیم؛ یک پایه‌ی غیر-SSL
بزرگ‌تر و بهتر تنظیم‌شده به کار آینده واگذار می‌شود پیش از اینکه محدودیت
۵ را در هر جهتی کاملاً حل‌شده تلقی کنیم.
\end{enumerate}

\textbf{اثر گسترده‌تر.} بهبود آشکارسازی دیپ‌فیک میان‌پیکره‌ای دومنظوره
است: همان سازوکار تطبیقی که به آشکارساز مدافع کمک می‌کند توزیع
درحال‌جابه‌جایی صدای واقعی استقرار را دنبال کند، می‌تواند، در اصل، به
مهاجمی بگوید کدام ویژگی‌های گفتار مصنوعی برای یک آشکارساز آسان‌تر است که
درست رتبه‌بندی کند، و پس کدام ویژگی‌ها برای فرار هدف قرار گیرند. ما این را
دلیلی برای پنهان‌داری روش نمی‌دانیم---پژوهش آشکارسازی روی گفتار مصنوعی
به‌طور گسترده منتشر می‌شود و آشکارسازها از پیش سمت ضعیف‌تر و کندتر این
مسابقه‌ی تسلیحاتی‌اند---اما آن را ملاحظه‌ای برای استقرار می‌دانیم: یک
اپراتور سامانه که به روش ما تکیه می‌کند باید تطبیق را فرایندی پیوسته در
برابر هدفی درحال‌جابه‌جایی ببیند، نه یک اصلاح یک‌باره، و نباید اعداد EER
این مقاله را تضمینی در برابر یک مهاجم باانگیزه با دسترسی تطبیقی به
آشکارساز بداند.

\section{نتیجه‌گیری}
آشکارسازی دیپ‌فیک صوتی میان‌پیکره‌ای نه به این دلیل شکست می‌خورد که
رتبه‌بندی از دست می‌رود، بلکه چون نقطه‌ی عملیاتی بدکالیبره است. با بهره‌بردن
از این، یک تطبیق بدون‌نظارت ساده در زمان آزمون---خودآموزی رتبه-مطمئن به‌علاوه‌ی
سازگاری کانال---تعمیم به پیکره‌ها و زبان‌های ندیده را در سه از چهار هدف
یک-پیکره-بیرون به‌طور قابل‌اعتماد بهبود می‌بخشد، جایی که TTA کمینه‌سازی
آنتروپی ناپایدار است، ستون‌فقرات غیر-SSL کاملاً فرومی‌پاشد، و یک پایه‌ی
خصمانه‌ی دامنه در زمان آموزش که تا همگرایی آموزش دیده با سه برابر محاسبات
عملکرد ضعیف‌تری دارد. روی هدف چهارم روش صادقانه خنثی است، و ما این را به
سیگنال رتبه‌بندی منبع ضعیف‌تر روی آن هدف نسبت می‌دهیم به‌جای پنهان کردن آن
با تنظیم به‌ازای هدف. یک تحلیل واگرایی دامنه‌ی جانشین نشان می‌دهد سازوکار
بازکالیبراسیون یک نمایش از پیش سازمان‌یافته بر اساس کلاس است، نه
تغییرناپذیری دامنه، که مستقلاً سازگار است با چرایی عملکرد ضعیف‌تر پایه‌ی
خصمانه‌ی دامنه با وجود بهینه‌سازی دقیقاً همان کمیتی که تحلیل ما نشان می‌دهد
گلوگاه نیست. یک تفکیک ۳۸‌زبانه بیشتر نشان می‌دهد استواری چندزبانه‌ی
آشکارساز به زبان‌های واقعاً ندیده در آموزش منبع هم گسترش می‌یابد، نه فقط
به زبان‌هایی که مصنوعات ترکیب‌سازی‌شان را دیده است.

\begin{latin}
\begin{thebibliography}{1}
\bibitem{wav2vec2aug} H.~Tak \emph{et al.}, ``Automatic speaker verification
spoofing and deepfake detection using wav2vec 2.0 and data augmentation,'' in
\emph{Proc.\ Odyssey}, 2022.
\bibitem{xlsrsls} Q.~Zhang \emph{et al.}, ``Audio deepfake detection with
self-supervised XLS-R and SLS classifier,'' 2024.
\bibitem{aasist} J.~Jung \emph{et al.}, ``AASIST: Audio anti-spoofing using
integrated spectro-temporal graph attention networks,'' in \emph{ICASSP}, 2022.
\bibitem{rawnet2} H.~Tak \emph{et al.}, ``End-to-end anti-spoofing with RawNet2,''
in \emph{ICASSP}, 2021.
\bibitem{asdg} Y.~Xie \emph{et al.}, ``Domain generalization via aggregation and
separation for audio deepfake detection,'' \emph{IEEE TIFS}, 2024.
\bibitem{crossdomain} ``Cross-domain audio deepfake detection: dataset and
analysis,'' 2024.
\bibitem{tent} D.~Wang \emph{et al.}, ``Tent: Fully test-time adaptation by
entropy minimization,'' in \emph{ICLR}, 2021.
\bibitem{dann} Y.~Ganin and V.~Lempitsky, ``Unsupervised domain adaptation by
backpropagation,'' in \emph{Proc.\ ICML}, 2015.
\bibitem{adabn} Y.~Li, N.~Wang, J.~Shi, J.~Liu, and X.~Hou, ``Revisiting batch
normalization for practical domain adaptation,'' \emph{arXiv:1603.04779}, 2016.
\bibitem{bendavid} S.~Ben-David, J.~Blitzer, K.~Crammer, A.~Kulesza,
F.~Pereira, and J.~W.~Vaughan, ``A theory of learning from different
domains,'' \emph{Machine Learning}, vol.\ 79, no.\ 1--2, pp.\ 151--175, 2010.
\bibitem{itw} N.~M\"uller, P.~Czempin, F.~Dieckmann, A.~Froghyar, and
K.~B\"ottinger, ``Does audio deepfake detection generalize?,'' in \emph{Proc.\
Interspeech}, 2022.
\end{thebibliography}
\end{latin}

\end{document}
